% Con número: \eqnAbsoluteHumidity
% Sin número: \eqnAbsoluteHumidity*
\NewDocumentCommand{\eqTotalMoles}{s}{
    \begin{equation}
        \IfBooleanTF{#1}{\nonumber}{\label{eq:TotalMoles}}
        \totalMole = \vaporMole + \gasMole
    \end{equation}
    \IfBooleanTF{#1}{}{Donde \vaporMole representa los moles del vapor y \gasMole los moles del gas (el componente que nos condensa).}
}

% y1 + y2 = 1
\NewDocumentCommand{\eqSumGaseosusFracMoles}{s}{
    \begin{equation}
    % Si hay asterisco (*), no numera. Si no, sí.
    \IfBooleanTF{#1}{\nonumber}{\label{eq:SumGaseosusFracMoles}}
    \vaporMoleFraction + \gasMoleFraction = 1
    \end{equation}
    % Solo muestra el texto si NO hay asterisco (la ecuación está numerada)
    \IfBooleanTF{#1}{}{Donde $\vaporMoleFraction$ representa la fracción molar del vapor y $\gasMoleFraction$ la fracción molar del gas.}
}

% y1 = 1 - y2
\NewDocumentCommand{\eqGasFracAsVaporFrac}{s}{
    \begin{subequations}
        \IfBooleanTF{#1}{\nonumber}{\label{eq:GasFracAsVaporFrac}} % Etiqueta para el grupo completo (opcional)
        \begin{align}
            \vaporMoleFraction + \gasMoleFraction = 1 \label{eq:SumVapGasFrac}\\ %No dejar espacio entre las ecuaciones
             \gasMoleFraction = 1 - \vaporMoleFraction \label{eq:GasFracAsVap}
        \end{align}
    \end{subequations}
}

\newcommand{\eqMoleVaporFractionAsAbsoluteHumidity}{
    \begin{equation}
        \label{eq:MoleVaporFractionAsAbsoluteHumidity}
        \vaporMoleFraction =  \frac{{\absoluteHumidity}}{{(\frac{\vaporMolecularWeight}{\gasMolecularWeight} + \absoluteHumidity)}}
    \end{equation}
}

\NewDocumentCommand{\eqMolecularWeightAverage}{s}{
    \begin{equation}
        \IfBooleanTF{#1}{\nonumber}{\label{eq:MolecularWeightAverage}}
        \averageMolecularWeight = (1 - \vaporMoleFraction) \gasMolecularWeight + \vaporMoleFraction \vaporMolecularWeight
    \end{equation}
    % \IfBooleanTF{#1}{}{Donde \vaporMole representa los moles del vapor y \gasMole los moles del gas (el componente que nos condensa).}
}